\documentclass[11pt]{article}
\usepackage{lmodern}
\usepackage[T1]{fontenc}
\usepackage{amsmath,amssymb,amsthm}
\usepackage[hidelinks]{hyperref}
\setlength{\emergencystretch}{1.5em}
\newtheorem{theorem}{Theorem}[section]
\newtheorem{lemma}[theorem]{Lemma}
\newtheorem{corollary}[theorem]{Corollary}
\newcommand{\Z}{\mathbf Z}
\newcommand{\Q}{\mathbf Q}
\newcommand{\C}{\mathbf C}
\newcommand{\ee}{\mathrm e}
\newcommand{\av}{\operatorname{av}}
\newcommand{\aug}{\operatorname{aug}}
\newcommand{\ord}{\operatorname{ord}}
\title{Recovering a central second derivative from primitive cyclotomic modular symbols}
\author{\small \href{mailto:admin@cosmosgeek.org}{admin@cosmosgeek.org}\qquad \href{mailto:b321014@iiit-bh.ac.in}{b321014@iiit-bh.ac.in}}
\date{}
\begin{document}
\maketitle

\begin{abstract}
For a rational elliptic newform of even root number we construct a real analytic test on its ordinary cyclotomic modular symbols. At conductor \(p^n\) this test pairs only with primitive multiplicative characters. The normalized pairings converge to an explicit multiple of the ordinary completed second derivative. A nonzero limit forces large complex pairings for finite antiderivatives in the cyclotomic group ring. For almost every fixed generator, the lower bound for an antiderivative of order \(r\) grows at least as the \((r-1)\)-st power of \(p^{n/2}/n^{2+\varepsilon}\). We give a separate example showing why divisibility by an augmentation power alone supplies no upper bound for these pairings.
\end{abstract}

\section{The central derivative and the cyclotomic test}

Let \(f(z)=\sum_{t\ge1}a_t\ee(tz)\) be the normalized weight-two newform of
an elliptic curve over \(\Q\) of conductor \(N\), with
\(\ee(x)=\exp(2\pi i x)\). Assume its root number is \(+1\).
Fix a good ordinary prime \(p\ge5\).
Write
\[
 \Lambda_f(s)=N^{s/2}(2\pi)^{-s}\Gamma(s)L(f,s),\qquad
 \Gamma_p=1+p\Z_p,\qquad \Delta=\mu_{p-1}\subset\Z_p^\times .
\]
For \(m=p^n\) put \(D_m=(\Z/m\Z)^\times\) and \(G_n=D_m/\Delta\).
The symbol \(\Delta\) in this quotient denotes its image modulo \(m\).
The real modular symbol used here is
\[
 S(x)=\Re S^{\mathrm c}(x),\qquad
 S^{\mathrm c}(x)=2\pi i\int_{i\infty}^{x} f(z)\,dz
\]
for a rational cusp \(x\), with the limiting integral along the vertical path.
The coefficients \(a_t\) are real, so \(S(-x)=S(x)\).
Let \(\alpha\) be the unit root of \(X^2-a_pX+p\), using a fixed
\(p\)-adic embedding of \(F=\Q(\alpha)\). Fix also a complex embedding for
every archimedean pairing below. Choose a positive real period
\(\Omega\) scaled by a rational factor so that the \(S(x)/\Omega\) are
rational and \(p\)-integral. Such a scaling follows from the modular
parametrization and the bounded denominators of its modular symbols.

The measure with the following convention belongs to
\(\mathcal O[[\Gamma_p]]\), where \(\mathcal O\) is the ring of integers
of the chosen \(p\)-adic completion of \(F\). Its finite projections have
coefficients in \(F\):
\begin{equation}\label{eq:measure}
 \mu_n(g)=\frac{\alpha^{-n}}{\Omega}
   \sum_{\substack{a\in D_m\\ a\mapsto g}}
       \left(S(a/m)-\alpha^{-1}S(pa/m)\right).
\end{equation}
Here and below the sum indexed by a residue class uses any integer
representative. The Hecke relation
\(\sum_{h=0}^{p-1}S((a+hm)/(pm))=a_pS(a/m)-S(pa/m)\)
gives compatibility under projection. This is the ordinary modular-symbol
construction with its tame character taken to be trivial
\cite{mtt}. Its augmentation is
\((1-\alpha^{-1})^2 L(f,1)/\Omega\).

Define real coefficients and a real analytic periodic function by
\begin{align}
 b_t&=2\int_1^\infty e^{-2\pi t y/\sqrt N}(\log y)^2\,dy,\nonumber\\
 B_j&=\sum_{h\ge0}a_{p^h} b_{jp^h} \qquad (p\nmid j),\nonumber\\
 K(x)&=2\sum_{\substack{j\ge1\\p\nmid j}}j B_j\cos(2\pi j x).
 \label{eq:kernel}
\end{align}
These series converge exponentially. For \(g\in G_n\) represented by
\(a\in D_m\) put
\[
 k_n(g)=\frac{1}{p-1}\sum_{\zeta\in\Delta}K((\zeta a\bmod m)/m).
\]
For \(x=\sum_{g\in G_n}x(g)[g]\in F[G_n]\), set
\begin{equation}\label{eq:pairing}
 R_n(x)=\frac{\Omega\alpha^n}{\varphi(m)}
            \sum_{g\in G_n} x(g)k_n(g).
\end{equation}
The fixed complex embedding applies to \(x(g)\) and \(\alpha^n\).
In particular this pairing is defined on the algebraic finite
projections, without choosing a map from \(\Q_p\) to \(\C\).

\begin{theorem}\label{thm:reconstruct}
For \(n\ge2\) the test \(k_n\) has zero sum on each fibre of
\(G_n\to G_{n-1}\). Consequently \(R_n\) uses only multiplicative
characters of conductor exactly \(p^n\). There is a
\(\delta=\delta(p)>0\) such that
\begin{equation}\label{eq:limit}
 R_n(\mu_n)=\frac{2}{p-1}\Lambda_f''(1)+O_{f,p}(m^{-\delta}).
\end{equation}
\end{theorem}

We prove the estimate directly. The relation between averages of
unprojected symbols and additive Fourier series has also been studied
by Diamantis, Hoffstein, K{\i}ral and Lee \cite{dhkl}.
The restriction to frequencies prime to \(p\) and the projection onto
\(G_n\) are both used in \eqref{eq:limit}.

\section{Sampling at additive cusps}

For a \(p\)-adic unit \(v\), reduce \(v\) modulo \(m\) and put
\[
 F_m(v)=\frac{1}{\varphi(m)}\sum_{a\in D_m}
                     S^{\mathrm c}(a/m)\ee(-av/m).
\]
For \(n\ge2\), let \(r_m(v)\) denote the least positive representative
of \(v\) modulo \(m/p\). Write \(c_m(h)=\sum_{a\in D_m}\ee(ah/m)\).

\begin{lemma}\label{lem:sampling}
Uniformly in \(v\in\Z_p^\times\), with \(X=m^{3/2}\) and \(r=r_m(v)\),
\[
 F_m(v)=\frac{c_m(r-v)}{\varphi(m)}\frac{a_r}{r}e^{-2\pi r/X}
           +O_{f,p,\varepsilon}(m^{-1/4+\varepsilon}).
\]
\end{lemma}

\begin{proof}
Split the integral defining \(S^{\mathrm c}(a/m)\) at height \(1/X\).
Choose \(d\) with \(Nad\equiv1\pmod m\) and put \(b=(Nad-1)/m\).
The matrix
\(\left(\begin{smallmatrix}Na&b\\ Nm&Nd\end{smallmatrix}\right)\)
lies in \(\Gamma_0(N)W_N\), with
\(W_N=\left(\begin{smallmatrix}0&-1\\N&0\end{smallmatrix}\right)\).
It sends \(-d/m+i t\) to \(a/m+i/(Nm^2t)\).
The weight-two transformation and root number \(+1\) give, with \(Y=Nm^2/X\),
\[
 S^{\mathrm c}(a/m)=
   \sum_{t\ge1}\frac{a_t}{t}\ee(ta/m)e^{-2\pi t/X}
   +\sum_{t\ge1}\frac{a_t}{t}\ee(-td/m)e^{-2\pi t/Y}.
\]
The normalized average of the first phase against \(\ee(-av/m)\) is
\(c_m(t-v)/\varphi(m)\). It is zero unless \(t\equiv v\pmod{m/p}\),
and has absolute value at most one. By
\(|a_t|\ll_{f,\varepsilon} t^{1/2+\varepsilon}\) \cite{deligne},
the terms \(t=r+h(m/p)\), \(h\ge1\), have total absolute value
\(O_{f,p,\varepsilon}(m^{-1/4+\varepsilon})\).

The average of the second phase is a Kloosterman sum divided by
\(\varphi(m)\), with first argument \(-v\) and second argument
\(-tN^{-1}\pmod m\). Its absolute value is \(O_p(m^{-1/2})\).
Indeed, put \(h=\lfloor n/2\rfloor\) and split the unit summation
variable as \(x+p^{n-h}z\), where \(x\) runs modulo \(p^{n-h}\)
and \(z\) modulo \(p^h\). Summation in \(z\) imposes
\(v x^2\equiv tN^{-1}\pmod{p^h}\). There are at most two unit roots
modulo \(p^h\), each with at most \(p\) lifts to \(p^{n-h}\).
The unnormalized bound is therefore at most \(2p^{n-h}\).
The dual contribution is \(O_{f,p,\varepsilon}(m^{-1/2}Y^{1/2+\varepsilon})\).
\end{proof}

\begin{proof}[Proof of Theorem \ref{thm:reconstruct}]
The expansion \eqref{eq:kernel} contains only \(p\)-unit frequencies.
Each has sum zero on the \(p\) lifts of a unit class modulo \(p^{n-1}\).
Averaging over \(\Delta\) preserves this property, which also removes
the lower-conductor term in \eqref{eq:measure}. Hence
\[
 R_n(\mu_n)=\frac{1}{\varphi(m)}\sum_{a\in D_m}S(a/m)k_n(a).
\]
The evenness and reality of \(k_n\) allow \(S\) to be replaced by
\(S^{\mathrm c}\) in this sum.

For a fixed positive integer \(j\) prime to \(p\), Lemma
\ref{lem:sampling} gives \(F_m(j)\to a_j/j\) and \(F_m(-j)\to0\).
For a Teichm\"uller root \(\zeta\ne\pm1\) of order \(d\), the
representative \(r=r_m(j\zeta)\) satisfies
\[
 m/p\ \mid\ j^{\varphi(d)}\Phi_d(r/j).
\]
This integer is nonzero, since a primitive root of order \(d>2\) is
irrational. Its absolute value is at most \((r+j)^{\varphi(d)}\).
The same argument applies to \(-\zeta\). Thus the terms
\(F_m(\pm j\zeta)\) tend to zero at a power rate, uniformly for
\(j\le C\log m\) with fixed \(C\). For \(\zeta=\pm1\) the estimates for
\(F_m(\pm j)\) are also uniform over this range, with the direct
exponential factor absorbed in the error. To discard larger \(j\),
use the exponential decay in \eqref{eq:kernel} and
\(\sup_{a\in D_m}|S^{\mathrm c}(a/m)|\ll_f 1+\log m\).
For the latter bound, the continued fraction of \(a/m\) gives
\(O(1+\log m)\) unimodular paths between consecutive convergents.
Integrals along these paths lie in a finite list up to sign, by
invariance of the weight-two differential and the finite index of
\(\Gamma_0(N)\). Choose \(C\) sufficiently large for the resulting
tail to have a power saving.

It follows that the limit is \(2\sum_{p\nmid j}a_j B_j/(p-1)\).
The root number gives
\(\Lambda_f''(1)=2\int_1^\infty f(iy/\sqrt N)(\log y)^2\,dy\)
by splitting the Mellin integral at one. Its Fourier expansion and
the identity \(a_{jp^h}=a_j a_{p^h}\) for \(p\nmid j\) give
\(\sum_{p\nmid j}a_j B_j=\Lambda_f''(1)\), as required.
\end{proof}

\section{Augmentation antiderivatives}

Let \(U\) be a topological generator of \(\Gamma_p\), also denoting
its image in \(G_n\) by \(U\). Assume \(L(f,1)=0\), so
\(\aug(\mu_n)=0\). For \(r\ge1\), define \(\eta_{n,r}\in F[G_n]\)
by
\[
 (1-[U])^r\eta_{n,r}=\mu_n,\qquad \aug(\eta_{n,r})=0.
\]
In characteristic zero, multiplication by \(1-[U]\) is invertible
on the kernel of augmentation. This proves uniqueness and existence.
Let \(I\) denote the augmentation ideal in the Iwasawa algebra.
If \(\mu\in I^{r+1}\), its quotient by \((1-[U])^r\) is integral
and has augmentation zero. By uniqueness, each finite specialization
agrees with the corresponding \(\eta_{n,r}\) in the chosen
\(p\)-adic embedding. In particular those specializations have
algebraic coefficients.

\begin{theorem}\label{thm:growth}
Fix \(\varepsilon>0\), and put
\(J_n=\lfloor p^{n/2}/n^{2+\varepsilon}\rfloor\).
For almost every fixed topological generator \(U\), with respect to
normalized Haar measure conditional on being a generator, the following
holds for all \(r\ge1\). If \(\Lambda_f''(1)\ne0\), then for all
sufficiently large \(n\),
\[
 \max_{|i|\le J_n+r}
      |R_n([U]^i\eta_{n,r})|
    \ \gg_{p,r}\ |\Lambda_f''(1)|\,J_n^{\,r-1}.
\]
\end{theorem}

\begin{proof}
Put \(A_{n,i}=R_n([U]^i\mu_n)\). After removal of the old term as
above, this is a linear combination of
\(F_m(\pm j\zeta U^i)\), with total absolute coefficient sum bounded
independently of \(n,i,U\). Additive Parseval and the continued-fraction
bound give
\[
 \sum_{v\bmod m}|F_m(v)|^2
   =\frac{m}{\varphi(m)^2}\sum_{a\in D_m}|S^{\mathrm c}(a/m)|^2
   \ll_{f,p}(1+\log m)^2,
\]
where \(F_m(v)\) is extended to all \(v\) by its defining sum.
For \(0<|i|\le J_n\), a given residue occurs under the map
\(U\mapsto j\zeta U^i\bmod m\) with probability
\(O_p(p^{v_p(i)}/m)\). Indeed, in a fixed generator coordinate on the
cyclic \(G_n\), this map multiplies a uniformly chosen unit exponent
by \(i\), followed by multiplication by the prescribed unit residue.
Cauchy--Schwarz yields
\[
 \mathbb E_U |A_{n,i}|
  \ll_{f,p} (1+\log m) p^{v_p(i)/2}/\sqrt m .
\]
The sum of \(p^{v_p(i)/2}\) over \(1\le i\le J\) is \(O_p(J)\).
The expected sum of \(|A_{n,i}|\) over \(0<|i|\le J_n\) is
therefore \(O_{f,p}(n^{-1-\varepsilon})\). These sums tend to
zero for almost every fixed \(U\).

Choose a smooth real function \(\psi\) supported in \((-1,1)\) with
\(\psi(0)=1\). By Theorem \ref{thm:reconstruct},
\[
 \sum_i \psi(i/J_n) A_{n,i}\ \longrightarrow\
          \frac{2}{p-1}\Lambda_f''(1).
\]
Set \(b_i'=R_n([U]^i\eta_{n,r})\). The relation defining \(\eta_{n,r}\)
gives \(A_{n,i}=(1-\mathcal S)^rb_i'\), where
\((\mathcal S b')_i=b_{i+1}'\). Summation by parts transfers these
differences to the sampled \(\psi\). The sum of absolute values of
its \(r\)-th finite differences is \(O_{\psi,r}(J_n^{1-r})\), by the
integral formula for a finite difference. Their support is contained
in \(|i|\le J_n+r\). This proves the claim.
\end{proof}

\begin{corollary}
Suppose \(\mu\in I^{r+1}\), and suppose that for a generator in the
full-measure set in Theorem \ref{thm:growth} the complex pairings of
the resulting integral quotient satisfy
\(\max_{|i|\le J_n+r}|R_n([U]^i\eta_{n,r})|=o(J_n^{r-1})\) on a
subsequence. Then \(L''(f,1)=0\).
\end{corollary}
\begin{proof}
The theorem gives \(\Lambda_f''(1)=0\). The even functional equation
and \(L(f,1)=0\) also give \(L'(f,1)=0\), so the ordinary and
completed second derivatives differ by a nonzero factor.
\end{proof}


For \(r=2\), three evaluations suffice for an exact signed limit.
Put \(B_{n,i}=R_n([U]^i\eta_{n,2})\). For a generator and window as in
Theorem \ref{thm:growth}, the second-difference identity gives
\begin{align*}
 \frac{B_{n,1+J_n}+B_{n,1-J_n}-2B_{n,1}}{J_n}
 &=\sum_{|i|<J_n}(1-|i|/J_n) A_{n,i}\\
 &\longrightarrow \frac{2}{p-1}\Lambda_f''(1).
\end{align*}
The finite equality follows by taking the second differences of the
triangular weights at their two endpoints and their midpoint. The
off-origin absolute-sum bound in the proof gives the limit. If
\(\mu\in I^3\), the three evaluations use its actual integral
second quotient. An \(o(J_n)\) upper bound on their displayed
numerator along a subsequence would imply ordinary second-derivative
vanishing. No such estimate is used here.

\section{Scope of the growth criterion}

For an explicit test of augmentation divisibility, take
\(u=1+p\sqrt{1+p}\in\Gamma_p\), choosing the \(p\)-adic square root
congruent to one modulo \(p\). This is a topological generator.
No positive power of \(u\) is rational: the two real embeddings of
\(u\) have different absolute values. Here \(1+p\) is nonsquare
for prime \(p\ge5\). For any fixed \(i\ne0\), all
\(\pm j\zeta u^i\), \(j\ge1\), are nonrational, as seen by raising
to a multiple of the order of \(\zeta\).
Their least positive representatives modulo \(p^{n-1}\) therefore
tend to infinity. Lemma \ref{lem:sampling} and exponential decay
imply
\[
 R_n([u]^i\mu_n)\ \longrightarrow\ 0 .
\]
For clarity, the sampling lemma bounds \(F_m(v)\) uniformly in both
\(m\) and \(p\)-unit \(v\), since
\(\sup_{r\ge1}|a_r|/r<\infty\). This justifies termwise passage to
the limit here. It follows by the binomial theorem that for every
fixed \(h\ge1\),
\[
 \lim_n R_n\big((1-[u])^h\mu_n\big)
    =\frac{2}{p-1}\Lambda_f''(1).
\]
The modified distributions are integral translates of the given
modular-symbol measure, combined with integer coefficients, and
belong to \(I^h\). For \(h\ge r+1\) their order-\(r\)
antiderivatives are themselves integral. Their shifted pairings
obey the same lower bound as in Theorem \ref{thm:growth} whenever
\(\Lambda_f''(1)\ne0\): each of the finitely many extra shifts
multiplies the frequencies in the expectation bound by a fixed unit,
so the proof applies with the same summable error. These examples show that
an archimedean upper estimate for the actual quotient in the corollary
requires further arithmetic input. No such estimate or general
rank conclusion is established here.


\section{What all fixed unit frequencies determine}

For \(h\in\Gamma_p\) define another test and its normalized pairing by
\begin{align*}
 k_{n,h}(a)&=\frac{2}{p-1}\sum_{\zeta\in\Delta}
        \cos\bigl(2\pi(\zeta ha\bmod m)/m\bigr),\\
 T_n(x;h)&=\frac{\Omega\alpha^n}{\varphi(m)}
                   \sum_{g\in G_n}x(g) k_{n,h}(g).
\end{align*}
For a positive rational \(q\) with \(v_p(q)=0\), let
\(\langle q\rangle=q/\omega(q)\in\Gamma_p\), where \(\omega\) is the
Teichm\"uller character. The map \(q\mapsto\langle q\rangle\) is
injective on these rationals, since the only positive rational root of
unity is one. Let \(H\) be its image.

\begin{theorem}\label{thm:profile}
For every fixed \(h\in\Gamma_p\) the limit \(T(h)=\lim_nT_n(\mu_n;h)\)
exists. Its value is
\[
 T(h)=
 \begin{cases}
 \dfrac{2a_j}{(p-1)j},&
   h=\langle j\rangle,\quad j\ge1,\quad p\nmid j,\\[3pt]
 0,&\text{otherwise}.
 \end{cases}
\]
Let \(\rho\) be a real signed Borel measure on \(\Gamma_p\) of finite
total variation, and write \(\rho_n\) for its pushforward to \(G_n\).
Use the fixed complex embedding on \(\mu_n\) in forming
\(\rho_n*\mu_n\). Then
\[
 \lim_{n\to\infty} T_n(\rho_n*\mu_n;h)=
                 \int_{\Gamma_p}T(hx)\,d\rho(x).
\]
If all cylinder masses of \(\rho\) belong to \(D^{-1}\Z\) for some
fixed positive integer \(D\), equality of this limiting profile with
\(T(h)\) for every \(h\) implies \(\rho=\delta_1\).
\end{theorem}

\begin{proof}
Each \(k_{n,h}\) has zero sum on the fibres down one level.
After removal of the lower-conductor term, the pairing is the average
of \(F_m(\pm\zeta h)\) over \(\zeta\), with both signs summed. For a
fixed \(p\)-adic unit \(v\), its least positive representatives modulo
\(p^{n-1}\) escape every bounded set unless \(v\) is a positive
integer. Lemma \ref{lem:sampling} thus gives the asserted profile.
The same lemma bounds all these sampled transforms uniformly, so
bounded convergence against the total variation of \(\rho\) proves
the convolution formula.

A real measure with masses in \(D^{-1}\Z\) has at most
\(D\|\rho\|_{\mathrm{TV}}\) nonzero cylinders in each level partition.
The number of nonzero cylinders is nondecreasing under refinement,
and hence eventually constant. Thereafter each nonzero cylinder has
one nonzero child and a zero cylinder has none. This proves that the
measure has finite support.

It remains to see that a nonzero finite sum
\(\sigma=\sum_i c_i\delta_{x_i}\) with distinct atoms cannot yield
the zero profile. In any coset of \(H\) meeting its support write the
atoms as \(x_i=x_0\langle q_i\rangle\). Choose the largest positive
rational \(q_*\) on that finite list and set
\(h=(x_0\langle q_*\rangle)^{-1}\). For every other atom in this
coset, \(hx_i=\langle q_i/q_*\rangle\) with \(0<q_i/q_*<1\).
Injectivity excludes every positive-integer contribution. Atoms in
other cosets also give zero. Thus \(\int T(hx)\,d\sigma(x)=c_*T(1)\ne0\).
Apply this statement to \(\sigma=\rho-\delta_1\).
\end{proof}

The fixed-denominator hypothesis has arithmetic content. Fix a
generator \(g\) of \(\Gamma_p\). Take independent digits
\(\epsilon_j=0,1\) with equal probabilities and let \(\nu\) be the
distribution of \(g^{\sum_{j\ge0}\epsilon_jp^j}\).
It is a nonatomic real probability measure. Its finite cylinder
masses are zero or \(2^{-(n-1)}\) on \(G_n\), so they also give
an element of \(\Z_p[[\Gamma_p]]\) of augmentation one. For \(s\ge1\)
each \(\nu^{*s}\) is nonatomic, whereas \(T\) is supported on a
countable set. Therefore Theorem \ref{thm:profile} gives, for every
\(h\in\Gamma_p\) and \(r\ge1\),
\[
 \lim_n T_n\big((\delta_1-\nu_n)^{*r}*\mu_n;h\big)=T(h).
\]
The smooth tests \eqref{eq:kernel} give the same conclusion term by
term, using their absolutely summable Fourier coefficients and the
uniform bound of Lemma \ref{lem:sampling}. These modified measures lie
in \(I^r\). The denominators of the auxiliary multiplier \(\nu_n\)
grow with \(n\), precisely the case excluded by the rigidity statement.
The rigidity applies to convolution multipliers under its stated
hypotheses. It supplies no growth bound for the integral
antiderivatives of a classical modular-symbol measure.


\section{Uniform tests at moving frequencies}

A version of the nonatomic example is uniformly invisible at all unit
frequencies, even when the frequency varies with the conductor. Use
the distribution \(\nu\) of \(g^{\sum_{j\ge0}\epsilon_jp^j}\) with
independent digits uniform on \(\{0,\ldots,p-2\}\). Its cylinder
masses at level \(n\) are zero or \((p-1)^{-(n-1)}\). In counting
measure,
\(\|\nu_n\|_2=(p-1)^{-(n-1)/2}\), while \(\|\nu_n\|_1=1\).
For \(s\ge1\) the same upper bound for \(\|\nu_n^{*s}\|_2\) follows
by convolution. For each fixed unit frequency and Teichm\"uller
factor, multiplication by that unit injects \(\Gamma_p/\Gamma_p^{p^{n-1}}\)
into \(D_m\). Additive Parseval from the proof of
Theorem \ref{thm:growth} and Cauchy--Schwarz therefore give
\[
 \sup_{h\in\Gamma_p}|T_n(\nu_n^{*s}*\mu_n;h)|
    \ll_{f,p}(1+\log m)(p-1)^{-(n-1)/2}.
\]
Thus the finite frequencies of \((\delta_1-\nu_n)^{*r}*\mu_n\) and
\(\mu_n\) differ uniformly by a power saving for every fixed \(r\).
The finite cylinder masses of the multiplier still have growing
denominators prime to \(p\).

In the opposite direction we recover individual cylinder increments
from the unit frequencies under a fixed-denominator hypothesis.
Define \(b'_d\) by \(L(f,s)^{-1}=\sum_{d\ge1}b'_d d^{-s}\)
in the half-plane of absolute convergence. The Euler factors and the
bound for \(a_d\) give \(|b'_d|\ll_{f,\varepsilon}d^{1/2+\varepsilon}\).
For \(T=\lfloor m^{1/4}\rfloor\) define
\[
 \mathcal D_n(x;h)=\frac12
   \sum_{\substack{d\le T\\p\nmid d}}\frac{b'_d}{d}
       T_n(x;h/\langle d\rangle).
\]
Let \({\bf1}_{n}(h)\) be the indicator that \(h\) is the identity
at level \(G_n\). Then, uniformly for \(h\in\Gamma_p\),
\begin{equation}\label{eq:spike}
 \mathcal D_n(\mu_n;h)=
   \frac{p}{(p-1)^2}\bigl({\bf1}_n(h)-p^{-1}{\bf1}_{n-1}(h)\bigr)
      +O_{f,p,\varepsilon}(m^{-1/8+\varepsilon}).
\end{equation}

To prove \eqref{eq:spike}, first apply the same sum over \(d\) to
\(F_m(v/d)\) for any \(v\in D_m\). Division here means multiplication
by the modular inverse. The errors in Lemma \ref{lem:sampling} sum
to \(O(m^{-1/4+\varepsilon}T^{1/2+\varepsilon})\). Put
\(h_d=r_m(v/d)\). The remaining terms with \(h_d>T^2\) are bounded
by \(O(T^{-1/2+\varepsilon})\), using the bounds for \(a_{h_d}\)
and \(b'_d\). For \(h_d\le T^2\), the integer \(t=d h_d\) lies below
\(m/p\) for all sufficiently large \(n\) and is congruent to \(v\) modulo
\(m/p\). Thus every such term has the same \(t\), if one exists.
The normalized Ramanujan factor is \(c_m(t-v)/\varphi(m)\),
independent of \(d\). Replacing the direct exponential by one costs
\(O((T^2/X)t^{-1/2+\varepsilon})\), since the sum of the absolute
values of \(b'_d a_{t/d}/t\) over divisors of \(t\) is
\(O_{f,\varepsilon}(t^{-1/2+\varepsilon})\).
If \(t>T\), that same bound gives a power saving.
If \(t\le T\), every divisor occurs and
\(\sum_{d\mid t}b'_d a_{t/d}={\bf1}_{t=1}\).
The raw result is therefore
\[
 \sum_{\substack{d\le T\\p\nmid d}}\frac{b'_d}{d}F_m(v/d)
 =
 \frac{p}{p-1}\left(
    {\bf1}_{v=1\bmod m}-p^{-1}{\bf1}_{v=1\bmod m/p}\right)
   + O_{f,p,\varepsilon}(m^{-1/8+\varepsilon}).
\]
Averaging over \(\zeta\in\Delta\) proves \eqref{eq:spike}. The factor
\(1/2\) cancels the two signs in the definition of \(T_n\). For
\(h\in\Gamma_p\) only the trivial Teichm\"uller factor can contribute
to the displayed raw indicators.

Let \(\rho_n\) now be compatible elements of \(\Q[G_n]\), with
\(\|\rho_n\|_1=m^{o(1)}\), and assume all their coefficients have a
fixed common denominator \(D\). Equation \eqref{eq:spike} gives
\begin{align*}
 \mathcal D_n(\rho_n*\mu_n;h)
 &=\frac{p}{(p-1)^2}
     \left(\rho_n(h^{-1})-p^{-1}\rho_{n-1}(\bar h^{-1})\right)\\
 &\quad+O_{f,p,\varepsilon}\left(m^{-1/8+\varepsilon}
                                      \|\rho_n\|_1\right).
\end{align*}
Here \(\bar h\) denotes reduction to \(G_{n-1}\). The cylinder
increments in parentheses belong to \((pD)^{-1}\Z\), so rounding
recovers them exactly at all sufficiently large levels.
Writing \(e_n(g)=\rho_n(g)-p^{-1}\rho_{n-1}(\bar g)\), one recovers
a mass along any chain of descendants by
\[
 \rho_t(g_t)=-\sum_{j\ge1}p^j e_{t+j}(g_{t+j})
       \qquad\text{in }\Q_p .
\]
The fixed denominator ensures that \(p^J\rho_{t+J}(g_{t+J})\)
tends to zero in \(\Q_p\) in the finite telescoping identity.
Projection then recovers the earlier levels.

In particular, suppose that for each \(\varepsilon>0\)
\[
 \sup_{h\in\Gamma_p}|T_n(\rho_n*\mu_n-\mu_n;h)|
                  \ll_\varepsilon m^{-1/4+\varepsilon}.
\]
Deconvolution costs \(O_{f,\varepsilon}(T^{1/2+\varepsilon})\).
Applying the increment formula to \(\rho_n-\delta_1\) shows that its
increments vanish identically at all sufficiently large levels,
and thus \(\rho_n=\delta_1\) at every level. This conclusion concerns
multipliers with the stated denominator and size bounds.

\begin{thebibliography}{9}
\bibitem{deligne}
P. Deligne, \emph{La conjecture de Weil. I},
Publ. Math. IH\'ES \textbf{43} (1974), 273--307.
\href{https://doi.org/10.1007/BF02684373}{doi:10.1007/BF02684373}.
\bibitem{dhkl}
N. Diamantis, J. Hoffstein, M. K{\i}ral and M. Lee,
\emph{Additive twists and a conjecture by Mazur, Rubin and Stein},
J. Number Theory \textbf{209} (2020), 1--36.
\href{https://doi.org/10.1016/j.jnt.2019.11.016}{doi:10.1016/j.jnt.2019.11.016}.
\bibitem{mtt}
B. Mazur, J. Tate and J. Teitelbaum,
\emph{On \(p\)-adic analogues of the conjectures of Birch and Swinnerton-Dyer},
Invent. Math. \textbf{84} (1986), 1--48.
\href{https://doi.org/10.1007/BF01388731}{doi:10.1007/BF01388731}.
\end{thebibliography}
\end{document}
